\section*{历史注记}
\phantomsection
\addcontentsline{toc}{section}{历史注记}

（方括号中的数字指本注记末尾的参考书目。）

乘法群 \(R_{+}^{*}\) （实数 \textgreater0 的）理论的历史，与 \textgreater0 的数的幂的概念的发展以及为表示它们所用的记法的发展密切相关。由同一个数的逐次幂组成的“几何数列”的想法可上溯到埃及人与巴比伦人；希腊数学家熟悉它，并且早在欧几里得{[}1{]}那里，我们就发现了等价于规则 \(a^{m}a^{n}=a^{m+n}\) （对整数指数 m, n\textgreater0）的一般陈述。在中世纪，法国数学家 N. 奥雷姆（14 世纪）重新发现了这条规则。他是第一个有 \textgreater0 的分数指数想法的人，其记法与我们的记法相似，并且以一般形式陈述了相应的计算规则；例如，他使用了我们今天应写作

\[
(a b) ^ {1 / n} = a ^ {1 / n} b ^ {1 / n}, \quad (a ^ {m}) ^ {p / q} = (a ^ {m p}) ^ {1 / q} [ 2 ].
\]

但奥雷姆的思想远远超前于他那个时代的数学，以致未能对其同时代人产生影响，他的论著不久便湮没无闻。一个世纪之后，N. 许凯重新陈述了欧几里得的规则；此外，他对他的方程中未知数的幂引入了指数记法，并且毫不迟疑地使用指数 o 与整指数 \(< 0\)（*）。这一次，尽管许凯的著作仍以手稿流传且似乎并未广泛传播，指数的“算术数列”与幂的“几何数列”之间的同构这一思想再也没有从视野中消失；它被施蒂费尔{[}4{]}推广到负指数与分数指数，并最终导致对数的定义和最初数张表的编制，这是由苏格兰人 J. 纳皮尔在 1614-1620 年{[}5{]}与瑞士人 J. 比尔吉（其著作直到 1620 年才发表，尽管他的思想可追溯到 17 世纪的最初几年）彼此独立地进行的。比尔吉在使用他的数表时借助于插值，隐含地假定了 R 与 \(R^{*}_{+}\) 之间所建立的同构的连续性；纳皮尔则相反，他在自己的定义中明确地表述了这一点（或者说，至少以当时流行的关于连续性的模糊概念所允许的那种明确程度）（*）。

在此我们的目的不是详述对数为数值计算所提供的服务；从理论观点看，对数的重要性始于无穷小演算的初创时期，即 \(\log(1+x)\) 与 \(e^{x}\) 的级数展开以及这些函数的微分性质的发现。就对数与指数函数的定义而言，直到 19 世纪中叶，人们一直直观地假定：对所有有理数 x 定义的函数 \(a^{x}\) 可以连续延拓到全体实数的集合上；直到实数概念被最终澄清并从有理数概念导出之后，人们才为这一连续延拓寻求严格的证明。一个类似的延拓原理，经过适当应用，再次成为 § 2 命题 1 与命题 2 的证明的基础；由它得到的不仅有指数函数与对数函数的定义，而且，正如我们将在第 VIII 章看到的，还有角的测度。

